
\documentclass[a4paper,10pt]{article}
\usepackage[utf8]{inputenc}
\usepackage[T1]{fontenc}
\usepackage[italian]{babel}
\usepackage{amsmath,amssymb,amsthm,amsfonts,mathtools,bm}
\usepackage{physics}
\usepackage{mathrsfs}
\usepackage{geometry}
\usepackage{array}
\usepackage{enumitem}
\usepackage{multicol}
\usepackage{xcolor}
\usepackage{tcolorbox}
\usepackage{titlesec}
\usepackage{hyperref}
\geometry{margin=1.4cm}
\setlength{\parindent}{0pt}
\setlength{\parskip}{2mm}
\definecolor{myblue}{RGB}{0,70,140}
\definecolor{mygray}{RGB}{240,240,240}
\titleformat{\section}{\large\bfseries\color{myblue}}{}{0em}{}
\titleformat{\subsection}{\normalsize\bfseries}{}{0em}{}
\newtcolorbox{formulaBox}{
colback=mygray,
colframe=myblue,
boxrule=0.6pt,
arc=2mm
}
\begin{document}
\begin{center}
{\LARGE \textbf{Formulario Completo — Metodi Matematici 1}}\\
\vspace{2mm}
Equazioni differenziali alle derivate parziali, serie di Fourier, operatori in spazi di Hilbert, completezza, operatori compatti.
\end{center}
\tableofcontents
\vspace{4mm}
%====================================================
\section{Serie di Fourier}
%====================================================
\subsection{Serie trigonometrica}
Per $f \in L^2(-\pi,\pi)$:
\[f(x) \sim \frac{a_0}{2}+\sum_{n=1}^{\infty}(a_n\cos nx+b_n\sin nx)\]
con
\[a_n=\frac{1}{\pi}\int_{-\pi}^{\pi}f(x)\cos(nx)dx\]
\[b_n=\frac{1}{\pi}\int_{-\pi}^{\pi}f(x)\sin(nx)dx\]
Forma complessa:
\[f(x)=\sum_{n\in\mathbb Z} c_n e^{inx}\]
\[c_n=\frac{1}{2\pi}\int_{0}^{2\pi}f(x)e^{-inx}dx\]
\subsection{Ortogonalità}
\[\int_{-\pi}^{\pi}e^{inx}e^{-imx}dx=2\pi\delta_{nm}\]
\[\int_0^{\pi}\sin(nx)\sin(mx)dx=\frac{\pi}{2}\delta_{nm}\]
\[\int_0^{\pi}\cos(nx)\cos(mx)dx=
\begin{cases}
\pi & n=m=0
\\
\frac{\pi}{2}\delta_{nm} & n,m>0
\end{cases}\]
\subsection{Convergenza}
\begin{itemize}
\item In $L^2$: sempre.
\item Puntuale nei punti di continuità.
\item Nei salti converge alla media:
\[\frac{f(x^+)+f(x^-)}{2}\]
\item Uniforme se $f$ è periodica e regolare ($C^1$ a tratti + continuità periodica).
\end{itemize}
\subsection{Parseval}
\[\frac{1}{\pi}\int_{-\pi}^{\pi}|f(x)|^2dx
=
\frac{a_0^2}{2}+\sum_{n=1}^{\infty}(a_n^2+b_n^2)\]
oppure
\[\frac{1}{2\pi}\int_{0}^{2\pi}|f(x)|^2dx
=
\sum_{n\in\mathbb Z}|c_n|^2\]
\subsection{Funzioni pari/dispari}
\begin{itemize}
\item $f$ pari $\Rightarrow$ solo coseni.
\item $f$ dispari $\Rightarrow$ solo seni.
\end{itemize}
\subsection{Sviluppi utili}
\[\sin^3x=\frac{3}{4}\sin x-\frac14\sin 3x\]
\[\cos^2x=\frac{1+\cos 2x}{2}\]
\[\sin^2x=\frac{1-\cos 2x}{2}\]
\[\sum_{n=1}^{\infty}\frac1{n^2}=\frac{\pi^2}{6}\]
\[\sum_{n=-\infty}^{\infty}\frac{e^{inx}}{n^2+a^2}
=\frac{\pi}{a\sinh(\pi a)}\cosh(a(\pi-x))\]
%====================================================
\section{Separazione delle Variabili}
%====================================================
Metodo standard:
\[u(x,t)=X(x)T(t)\]
Sostituire nella PDE e separare:
\[\frac{T'}{T}=\frac{X''}{X}=-\lambda\]
Si ottengono due ODE.
\subsection{Caso classico}
\[X''+\lambda X=0\]
Condizioni di Dirichlet:
\[X(0)=X(\pi)=0\]
Autovalori:
\[\lambda_n=n^2\]
Autofunzioni:
\[X_n(x)=\sin(nx)\]
Condizioni di Neumann:
\[X'(0)=X'(\pi)=0\]
Autofunzioni:
\[X_n(x)=\cos(nx)\]
%====================================================
\section{Equazione del Calore}
%====================================================
\subsection{Forma standard}
\[\partial_t u=\partial_{xx}u\]
con $u(0,t)=u(\pi,t)=0$.
\subsection{Soluzione generale}
\[u(x,t)=\sum_{n=1}^{\infty} a_n e^{-n^2 t}\sin(nx)\]
con
\[a_n=\frac{2}{\pi}\int_0^{\pi}f(x)\sin(nx)dx\]
\subsection{Metodo generale}
\begin{enumerate}
\item Separare variabili.
\item Risolvere problema agli autovalori.
\item Sviluppare dato iniziale in serie di autofunzioni.
\item Sommare i modi.
\end{enumerate}
\subsection{Equazione con termine convettivo}
Per
\[\partial_tu=\partial_{xx}u+\partial_xu+u\]
si usa
\[u=e^{\alpha t+\beta x}v\]
scegliendo $\alpha,\beta$ per eliminare termini di primo ordine.
Tipicamente:
\[\beta=-\frac12\]
\[\alpha=-\frac34\]
in modo da ottenere:
\[\partial_t v=\partial_{xx}v\]
%====================================================
\section{Equazione della Corda Vibrante}
%====================================================
\[\partial_{tt}u=\partial_{xx}u\]
con estremi fissi:
\[u(0,t)=u(\pi,t)=0\]
\subsection{Soluzione generale}
\[u(x,t)=\sum_{n=1}^{\infty}
\left[
A_n\cos(nt)+B_n\sin(nt)
\right]\sin(nx)\]
Coefficienti:
\[A_n=\frac2\pi\int_0^{\pi}u(x,0)\sin(nx)dx\]
\[B_n=\frac2{n\pi}\int_0^{\pi}u_t(x,0)\sin(nx)dx\]
\subsection{Energia}
\[E(t)=\int_0^{\pi}
\left[
(u_t)^2+(u_x)^2
\right]dx\]
Per soluzioni regolari:
\[\frac{dE}{dt}=0\]
quindi energia conservata.
\subsection{Formula di d'Alembert}
Su $\mathbb R$:
\[u(x,t)=F(x-t)+G(x+t)\]
oppure
\[u(x,t)=\frac{f(x-t)+f(x+t)}{2}
+\frac12\int_{x-t}^{x+t}g(s)ds\]
%====================================================
\section{Equazione di Laplace}
%====================================================
\[\Delta u=u_{xx}+u_{yy}=0\]
\subsection{Rettangolo con Dirichlet}
In $0<x<\pi$, $0<y<\pi$:
\[u(x,y)=\sum_{n=1}^{\infty}
\left[
A_n\sinh(ny)+B_n\sinh(n(\pi-y))
\right]\sin(nx)\]
\subsection{Metodo pratico}
Per condizioni su un solo lato:
\[u(x,y)=\sum a_n\sin(nx)\frac{\sinh(n(\pi-y))}{\sinh(n\pi)}\]
\subsection{Simmetrie utili}
Se il problema è invariante per:
\[(x,y)\leftrightarrow(y,x)\]
allora:
\[u(x,y)=u(y,x)\]
sulla diagonale:
\[u_x=u_y\]
quindi derivata normale nulla.
Se simmetria rispetto a:
\[x+y=\pi\]
allora:
\[u_x+u_y=0\]
sulla diagonale.
%====================================================
\section{Coordinate Polari}
%====================================================
\[\Delta u=u_{rr}+\frac1r u_r+\frac1{r^2}u_{\phi\phi}\]
Ansatz:
\[u(r,\phi)=R(r)\Phi(\phi)\]
Equazioni separate:
\[\Phi''+\lambda\Phi=0\]
\[r^2R''+rR'-\lambda R=0\]
\subsection{Soluzioni radiali}
Per $\lambda=n^2$:
\[R(r)=Ar^n+Br^{-n}\]
Limitatezza in $r=0$:
\[B=0\]
\subsection{Serie armoniche}
\[u(r,\phi)=a_0+\sum_{n=1}^{\infty}r^n
(a_n\cos n\phi+b_n\sin n\phi)\]
%====================================================
\section{Spazi di Hilbert}
%====================================================
\subsection{Definizioni}
Prodotto scalare:
\[(f,g)=\int f\bar g\]
Norma:
\[\|f\|=\sqrt{(f,f)}\]
Ortogonalità:
\[(f,g)=0\]
Base ortonormale:
\[(e_n,e_m)=\delta_{nm}\]
Sviluppo:
\[f=\sum_n(f,e_n)e_n\]
Bessel:
\[\sum|(f,e_n)|^2\le \|f\|^2\]
Parseval:
\[\sum|(f,e_n)|^2=\|f\|^2\]
%====================================================
\section{Completezza}
%====================================================
\subsection{Criterio fondamentale}
$\{f_n\}$ completo se:
\[(f_n,g)=0\ \forall n
\Rightarrow g=0\]
\subsection{Fatti utili}
\begin{itemize}
\item I polinomi sono completi in $L^2(0,1)$.
\item Gli esponenziali $e^{inx}$ sono completi.
\item Seni e coseni formano basi complete.
\item Se una combinazione lineare permette di ottenere una base completa, allora l'insieme è completo.
\end{itemize}
\subsection{Tecniche tipiche d'esame}
\begin{itemize}
\item Usare convergenza dominata.
\item Integrare per parti.
\item Usare Fubini.
\item Ricondursi a Fourier.
\item Usare parità/disparità.
\end{itemize}
%====================================================
\section{Operatori Lineari}
%====================================================
\subsection{Definizione}
\[T(\alpha x+\beta y)=\alpha Tx+\beta Ty\]
\subsection{Operatore limitato}
\[\|Tf\|\le C\|f\|\]
Norma:
\[\|T\|=\sup_{\|f\|=1}\|Tf\|\]
\subsection{Aggiunto}
\[(Tf,g)=(f,T^{\dagger}g)\]
\subsection{Autoaggiunto}
\[T=T^{\dagger}\]
Autovalori reali.
\subsection{Unitario}
\[T^{\dagger}T=I\]
Conserva la norma.
\subsection{Normale}
\[TT^{\dagger}=T^{\dagger}T\]
Autovettori relativi ad autovalori distinti ortogonali.
%====================================================
\section{Operatori Integrali}
%====================================================
Forma:
\[Tf(x)=\int K(x,y)f(y)dy\]
\subsection{Compattezza}
Se
\[K\in L^2\]
allora $T$ è compatto (Hilbert-Schmidt).
\subsection{Aggiunto}
\[T^{\dagger}f(y)=\int \overline{K(x,y)}f(x)dx\]
%====================================================
\section{Operatori Diagonali}
%====================================================
Se
\[Te_n=a_ne_n\]
allora:
\[\|T\|=\sup_n|a_n|\]
\[T^{\dagger}e_n=\bar a_ne_n\]
Autovalori:
\[\lambda_n=a_n\]
Compatto se:
\[a_n\to0\]
%====================================================
\section{Operatori di Riflessione e Simmetria}
%====================================================
Operatori tipici:
\[Tf(x)=f(-x)\]
\[Tf(x)=\operatorname{sgn}(x)f(x)\]
\subsection{Funzioni pari/dispari}
\[f=f_P+f_D\]
con
\[f_P(x)=\frac{f(x)+f(-x)}2\]
\[f_D(x)=\frac{f(x)-f(-x)}2\]
Proiettore sulle dispari:
\[Df=f_D\]
%====================================================
\section{Autovalori e Autovettori}
%====================================================
\[Tf=\lambda f\]
\subsection{Tecniche pratiche}
\begin{itemize}
\item Provare funzioni pari/dispari.
\item Provare autofunzioni della base.
\item Calcolare $T^2$.
\item Usare decomposizione ortogonale.
\end{itemize}
\subsection{Fatti importanti}
Se $T$ compatto autoaggiunto:
\begin{itemize}
\item autovalori reali;
\item autovettori ortogonali;
\item base ortonormale di autovettori.
\end{itemize}
%====================================================
\section{Convergenze di Operatori}
%====================================================
\subsection{Convergenza forte}
\[T_nx\to Tx\]
per ogni $x$.
\subsection{Convergenza in norma}
\[\|T_n-T\|\to0\]
Più forte della convergenza forte.
\subsection{Importante}
\[\|T_n-T\|\to0
\Rightarrow T_nx\to Tx\]
non viceversa.
%====================================================
\section{Identità Utilissime da Esame}
%====================================================
\[\sin(a\pm b)=\sin a\cos b\pm\cos a\sin b\]
\[\cos(a\pm b)=\cos a\cos b\mp\sin a\sin b\]
\[\sin a\sin b=\frac12[\cos(a-b)-\cos(a+b)]\]
\[\cos a\cos b=\frac12[\cos(a-b)+\cos(a+b)]\]
\[\sin a\cos b=\frac12[\sin(a+b)+\sin(a-b)]\]
\[\cosh x=\frac{e^x+e^{-x}}2\]
\[\sinh x=\frac{e^x-e^{-x}}2\]
\[\cos(ix)=\cosh x\]
\[\sin(ix)=i\sinh x\]
%====================================================
\section{Teoremi Utili da Citare}
%====================================================
\subsection{Convergenza dominata}
Se:
\[f_n\to f\]
q.o. e
\[|f_n|\le g\in L^1\]
allora
\[\int f_n\to\int f\]
\subsection{Fubini}
Permette di scambiare ordine di integrazione.
\subsection{Riemann-Lebesgue}
Se $f\in L^1$:
\[\int f(x)e^{inx}dx\to0\]
\subsection{Weierstrass}
I polinomi sono densi in $C([a,b])$.
\subsection{Sturm-Liouville}
Problemi autoaggiunti:
\[-(py')'+qy=\lambda wy\]
producono autofunzioni ortogonali rispetto a:
\[(f,g)_w=\int f\bar g w\]
%====================================================
\section{Schema Risolutivo Tipico PDE}
%====================================================
\begin{enumerate}
\item Identificare PDE.
\item Scrivere condizioni al bordo.
\item Applicare separazione variabili.
\item Risolvere problema agli autovalori.
\item Sviluppare dato iniziale/contorno.
\item Determinare coefficienti con ortogonalità.
\item Verificare simmetrie.
\item Controllare convergenza.
\end{enumerate}
%====================================================
\section{Schema Risolutivo Tipico Operatori}
%====================================================
\begin{enumerate}
\item Verificare linearità.
\item Stimare norma.
\item Calcolare aggiunto.
\item Verificare:
\begin{itemize}
\item autoaggiunto;
\item unitario;
\item normale;
\item compatto.
\end{itemize}
\item Trovare autovalori.
\item Studiare convergenze.
\end{enumerate}
%====================================================
\section{Trucchi Frequenti negli Esami}
%====================================================
\begin{itemize}
\item Usare simmetria per annullare derivate normali.
\item Usare decomposizione pari/dispari.
\item Calcolare prima $T^2$.
\item Trasformare PDE in equazione standard con cambio esponenziale.
\item Usare completezza tramite Fourier.
\item Per operatori integrali usare kernel.
\item Per compattezza usare Hilbert-Schmidt.
\item Per operatori diagonali usare direttamente la base.
\item Per serie di Fourier usare identità trigonometriche.
\item Per convergenza forte verificare punto per punto sui coefficienti.
\end{itemize}

%====================================================
\section{Serie Numeriche Utili}
%====================================================
\subsection{Serie geometriche}
\[\sum_{n=0}^{\infty} q^n = \frac{1}{1-q}
\qquad |q|<1\]
\[\sum_{n=1}^{\infty} nq^{n-1}=\frac{1}{(1-q)^2}
\qquad |q|<1\]
\[\sum_{n=1}^{\infty} nq^n=\frac{q}{(1-q)^2}\]
\subsection{Serie armoniche e p-serie}
\[\sum_{n=1}^{\infty}\frac1n\]
diverge.
\[\sum_{n=1}^{\infty}\frac1{n^p}\]
converge se e solo se:
\[p>1\]
\subsection{Serie alternate}
Criterio di Leibniz:
Se:
\begin{itemize}
\item $a_n\ge0$
\item $a_n$ decrescente
\item $a_n\to0$
\end{itemize}
allora:
\[\sum (-1)^n a_n\]
converge.
\subsection{Serie di Taylor fondamentali}
\[e^x=\sum_{n=0}^{\infty}\frac{x^n}{n!}\]
\[\sin x=\sum_{n=0}^{\infty}(-1)^n\frac{x^{2n+1}}{(2n+1)!}\]
\[\cos x=\sum_{n=0}^{\infty}(-1)^n\frac{x^{2n}}{(2n)!}\]
\[\ln(1+x)=\sum_{n=1}^{\infty}(-1)^{n+1}\frac{x^n}{n}
\qquad |x|<1\]
\[\frac1{1-x}=\sum_{n=0}^{\infty}x^n
\qquad |x|<1\]
%====================================================
\section{Tecniche Generali di Risoluzione}
%====================================================
\subsection{Strategia generale per PDE}
Per quasi tutti gli esercizi del corso il metodo principale è la separazione delle variabili.
Bisogna:
\begin{enumerate}
\item identificare il tipo di PDE;
\item guardare dominio e bordo;
\item imporre condizioni iniziali e al contorno;
\item scegliere una forma separabile;
\item risolvere il problema agli autovalori;
\item sviluppare il dato iniziale nella base di autofunzioni.
\end{enumerate}
Nella pratica si pone:
\[u(x,t)=X(x)T(t)\]
oppure:
\[u(x,y)=X(x)Y(y)\]
Si sostituisce nella PDE e si divide per il prodotto delle funzioni ottenendo due membri dipendenti da variabili diverse.
Poiché devono coincidere per ogni variabile, entrambi devono essere costanti.
\subsection{Scelta della costante di separazione}
Di solito si pone:
\[X''+\lambda X=0\]
per ottenere seni e coseni.
I tre casi fondamentali sono:
\begin{itemize}
\item $\lambda<0$: esponenziali/iperboliche;
\item $\lambda=0$: soluzioni lineari;
\item $\lambda>0$: soluzioni oscillanti.
\end{itemize}
Con Dirichlet o Neumann spesso l'unico caso non banale è $\lambda>0$.
\subsection{Uso dell'ortogonalità}
Dopo aver trovato le autofunzioni si espande il dato iniziale:
\[f(x)=\sum a_n\sin(nx)\]
oppure:
\[f(x)=\sum a_n\cos(nx)\]
Per trovare i coefficienti si moltiplica per l'autofunzione corrispondente e si integra.
L'ortogonalità elimina tutti i termini tranne uno.
Questo è il meccanismo fondamentale dietro Fourier.
\subsection{Condizioni non omogenee}
Se il bordo non è omogeneo conviene introdurre una soluzione stazionaria.
Si pone:
\[u=v+w\]
con $w$ scelta in modo da soddisfare le condizioni al bordo.
La funzione $v$ soddisferà condizioni omogenee e potrà essere sviluppata in serie.
\subsection{Operatori lineari: schema mentale}
Quando compare un operatore conviene controllare:
\begin{enumerate}
\item linearità;
\item limitatezza;
\item aggiunto;
\item autoaggiuntezza;
\item unitarietà;
\item normalità;
\item compattezza;
\item autovalori/autovettori.
\end{enumerate}
\subsection{Tecniche tipiche per autovalori}
Tecniche frequenti:
\begin{itemize}
\item usare funzioni pari/dispari;
\item usare Fourier;
\item usare esponenziali complessi;
\item provare direttamente le autofunzioni della base.
\end{itemize}
\subsection{Completezza}
Per mostrare completezza:
\begin{enumerate}
\item supporre ortogonalità a tutto il sistema;
\item dimostrare che tutti i coefficienti di Fourier sono nulli;
\item dedurre che la funzione è nulla.
\end{enumerate}
\subsection{Trucchi utili}
\begin{itemize}
\item usare parità/disparità per semplificare integrali;
\item usare identità trigonometriche prodotto-somma;
\item usare simmetrie geometriche;
\item integrare per parti nelle serie di Fourier;
\item usare convergenza dominata e Fubini.
\end{itemize}
\end{document}
```
